\documentstyle[11pt]{article}



%\renewcommand{\baselinestretch}{1.1}

\thispagestyle{empty}

\begin{document}



\begin{center}

{\large \bf NEW WAYS OF DEFINING}

\end{center}

\begin{center}

{\large \bf FILIFORM LIE ALGEBRAS}

\end{center}



\vspace{0.15cm}



\begin{center}

{\bf J.C. Benjumea, F.J. Echarte, D. Fern\'andez,}

\end{center}

\begin{center}

{\bf M.C. M\'arquez, J. N\'u{\~n}ez and F.
Ram\'{\i}rez}

\end{center}



\vspace{0.1cm}



\begin{center}

{\small Dept. Geometr\'{\i}a y Topolog\'{\i}a. Facultad de Matem�ticas. \\

%\end{center}

%\begin{center}

Universidad de Sevilla.

Aptdo 1160. 41080-Sevilla (Espa�a). }\\

\end{center}



\vspace{0.15cm}



\begin{tabular}{lllll}

E-mails & JCB: & acevedo@cica.es &DF: & desamfer@cica.es \\

& FJE: & echarte@cica.es & MCM: & cmgarcia@cica.es \\

& JN: &jnvaldes@cica.es 

& FR: & frlopez@cica.es 

\end{tabular}



\vspace{0.3cm}



\begin{center} 

{\large \bf Extended Abstract}

\end{center}

\vspace{0.20cm}



As it is already known, 
filiform Lie algebras, which are a subset of nilpotent Lie
algebras, were defined 
by Vergne in $1966.$ 



At present, the traditional way of defining complex Lie algebras
in ge\-ne\-ral, and complex filiform 
Lie algebras in particular, consists in giving all the non-null
brackets between elements of a basis of the
algebra. 



Then, the main goal of this paper is to show 
two new ways of de\-fi\-ning 
complex filiform Lie algebras, in which the number of
non-null  brackets needed is reduced.
To do this, we use the invariants $i$ and $j$ of complex filiform
Lie  algebras, previously defined by some of the
authors of this paper in 
earlier papers. In fact, we prove that a 
complex filiform Lie algebra of dimension $n$ is completely
defined  with respect to an adapted basis $\{e_1,
e_2, \ldots, e_n\}$ 
if either only the brackets
$[e_h, e_n] \quad (i \leq h < n)$ or only the
brackets
$[e_k, e_{k+1}] \quad (j \leq k < n)$ are known, where $i$ and
$j$ are the invariants previosly mentioned.



We also give in this paper an algorithmic
procedure  which allows to compute all the
non-null brackets appearing in
traditional definition starting of the only ones appearing in
ours.  This algorithm involves lots of
computations, which can be 
made by using any symbolic computation package. 





\end{document}

